\documentclass[10pt, letterpaper]{article}
% Packages:
\usepackage[
    ignoreheadfoot, % set margins without considering header and footer
    top=2 cm, % seperation between body and page edge from the top
    bottom=2 cm, % seperation between body and page edge from the bottom
    left=2 cm, % seperation between body and page edge from the left
    right=2 cm, % seperation between body and page edge from the right
    footskip=1.0 cm, % seperation between body and footer
    % showframe % for debugging 
]{geometry} % for adjusting page geometry
\usepackage{titlesec} % for customizing section titles
\usepackage{tabularx} % for making tables with fixed width columns
\usepackage{array} % tabularx requires this
\usepackage[dvipsnames]{xcolor} % for coloring text
\definecolor{primaryColor}{RGB}{0, 0, 0} % define primary color
\usepackage{enumitem} % for customizing lists
\usepackage{fontawesome5} % for using icons
\usepackage{amsmath} % for math
\usepackage[
    pdftitle={Adib Akkari's CV},
    pdfauthor={Adib Akkari},
    pdfcreator={LaTeX with RenderCV},
    colorlinks=true,
    urlcolor=primaryColor,
    draft
]{hyperref} % for links, metadata and bookmarks
\usepackage[pscoord]{eso-pic} % for floating text on the page
\usepackage{calc} % for calculating lengths
\usepackage{bookmark} % for bookmarks
\usepackage{lastpage} % for getting the total number of pages
\usepackage{changepage} % for one column entries (adjustwidth environment)
\usepackage{paracol} % for two and three column entries
\usepackage{ifthen} % for conditional statements
\usepackage{needspace} % for avoiding page brake right after the section title
\usepackage{iftex} % check if engine is pdflatex, xetex or luatex
% Ensure that generate pdf is machine readable/ATS parsable:
\ifPDFTeX
    \input{glyphtounicode}
    \pdfgentounicode=1
    \usepackage[T1]{fontenc}
    \usepackage[utf8]{inputenc}
    \usepackage{lmodern}
\fi
\usepackage{roboto}
% Some settings:
\raggedright
\AtBeginEnvironment{adjustwidth}{\partopsep0pt} % remove space before adjustwidth environment
\pagestyle{empty} % no header or footer
\setcounter{secnumdepth}{0} % no section numbering
\setlength{\parindent}{0pt} % no indentation
\setlength{\topskip}{0pt} % no top skip
\setlength{\columnsep}{0.15cm} % set column seperation
\pagenumbering{gobble} % no page numbering
\titleformat{\section}{\needspace{4\baselineskip}\bfseries\large}{}{0pt}{}[\vspace{1pt}\titlerule]
\titlespacing{\section}{
    % left space:
    -1pt
}{
    % top space:
    0.3 cm
}{
    % bottom space:
    0.2 cm
} % section title spacing
\renewcommand\labelitemi{$\vcenter{\hbox{\small$\bullet$}}$} % custom bullet points
\newenvironment{highlights}{
    \begin{itemize}[
        topsep=0.10 cm,
        parsep=0.10 cm,
        partopsep=0pt,
        itemsep=0pt,
        leftmargin=0 cm + 10pt
    ]
}{
    \end{itemize}
} % new environment for highlights
\newenvironment{highlightsforbulletentries}{
    \begin{itemize}[
        topsep=0.10 cm,
        parsep=0.10 cm,
        partopsep=0pt,
        itemsep=0pt,
        leftmargin=10pt
    ]
}{
    \end{itemize}
} % new environment for highlights for bullet entries
\newenvironment{onecolentry}{
    \begin{adjustwidth}{
        0 cm + 0.00001 cm
    }{
        0 cm + 0.00001 cm
    }
}{
    \end{adjustwidth}
} % new environment for one column entries
\newenvironment{twocolentry}[2][]{
    \onecolentry
    \def\secondColumn{#2}
    \setcolumnwidth{\fill, 4.5 cm}
    \begin{paracol}{2}
}{
    \switchcolumn \raggedleft \secondColumn
    \end{paracol}
    \endonecolentry
} % new environment for two column entries
\newenvironment{threecolentry}[3][]{
    \onecolentry
    \def\thirdColumn{#3}
    \setcolumnwidth{, \fill, 4.5 cm}
    \begin{paracol}{3}
    {\raggedright #2} \switchcolumn
}{
    \switchcolumn \raggedleft \thirdColumn
    \end{paracol}
    \endonecolentry
} % new environment for three column entries
\newenvironment{header}{
    \setlength{\topsep}{0pt}\par\kern\topsep\centering\linespread{1.5}
}{
    \par\kern\topsep
} % new environment for the header
\newcommand{\placelastupdatedtext}{% \placetextbox{<horizontal pos>}{<vertical pos>}{<stuff>}
  \AddToShipoutPictureFG*{% Add <stuff> to current page foreground
    \put(
        \LenToUnit{\paperwidth-2 cm-0 cm+0.05cm},
        \LenToUnit{\paperheight-1.0 cm}
    ){\vtop{{\null}\makebox[0pt][c]{
        \small\color{gray}\textit{Last updated in July 2025}\hspace{\widthof{Last updated in July 2025}}
    }}}%
  }%
}%
% save the original href command in a new command:
% new command for external links:
\begin{document}

\begin{center}
    \textbf{Adib Akkari} · adibakkari@gmail.com · (514) 984 9929 · Montreal, QC \\
    French (Delf B1) \& English · LinkedIn: linkedin.com/in/adib-akkari · GitHub: github.com/adssib
\end{center}
% \section{Introduction}
% \begin{onecolentry}
% Software Engineering student at Concordia University with a focus on data engineering, cloud computing, and applied artificial intelligence. Skilled in Python, cloud platforms, and frameworks for machine learning and software development.\end{onecolentry}
\vspace{-0.6cm}

\section{Education}        
\begin{twocolentry}{
    September 2022 - April 2027
}
    Concordia University, Bachelor in Software Engineering  \\
    % {\small \textbf{Relevant Coursework:} Operating Systems, Computer Networks, Database Systems, Intro to Deep Learning \& Machine Learning}
\end{twocolentry}

\vspace{-0.2cm}

\section{Experience}
\begin{twocolentry}{
    January 2026 – August 2026
}
\textbf{Software Developer Internship}, \textit{Ericsson} – Montreal, QC
\end{twocolentry}
\begin{onecolentry}
\textit{DevOps \& Infrastructure}
\begin{itemize}
    \item Owned daily reliability of 20+ release-critical CI/CD pipelines, driving cross-team root-cause analysis, eliminating false-positive runs through pipeline enhancements, and documenting failure patterns to reduce resolution time
    \item Reduced performance and stability test report generation time by up to 33\% by replacing Selenium-based UI extraction with direct Grafana API calls, improving pipeline reliability across version upgrades
\end{itemize}
\textit{Software Development \& Testing}
\begin{itemize}
    \item Developed features across a cloud-native IAM platform handling authentication for 130M+ users, with unit, component, performance and functional test coverage
    \item Benchmarked 15+ OAuth APIs and Kafka event pipelines using JMeter with controlled pod isolation (HPA limiting, synthetic data, scaled dependencies)
    \item Developed AI agent skills automating performance testing, documentation, report generation, and knowledge-sharing workflows across QA/LSV teams
\end{itemize}
\end{onecolentry}

\vspace{5 pt - 0.0 cm}
\begin{twocolentry}{
    January 2024 – January 2026
}
\textbf{Software Developer} \textbf{Internship}, \textit{Novatek International} – Montreal, QC
\end{twocolentry}
\vspace{0.10 cm}
\begin{onecolentry}
    \begin{itemize}
        \item Led end-to-end planning and delivery of a mission-critical VCS migration involving IT, QA, and senior VPs, rescuing 20+ years of source code and 200,000+ files from unstable legacy infrastructure
        \item Engineered an internal migration tool transferring 700,000+ commits with full history, tags, and access controls by reverse-engineering an undocumented legacy VCS system
        \item Deployed VisualSVN infrastructure integrated into CI/CD pipelines, then built an SVN-to-GitLab CLI tool enabling the company's transition from legacy VCS to Git
    \end{itemize}
\end{onecolentry}

\section{Projects}

\begin{twocolentry}{
    April 2026
}
\textbf{Natural Language to SQL Engine} | Python, PyTorch, Flask, Next.js
  \end{twocolentry}
  \vspace{0.10 cm}
  \begin{onecolentry}
      \begin{highlights}
  \item Built an end-to-end data pipeline unifying 134K examples across 5 sources, then fine-tuned Llama 3 8B with QLoRA (4-bit NF4, LoRA rank 16) for SQL dialect generation          
  \item Designed and deployed AskDB, a full-stack web application where users provide a DDL schema and query their database in natural english language
      \end{highlights}
  \end{onecolentry}
  \vspace{0.2cm}


% \begin{twocolentry}{
%     November 2025
% }
% \textbf{QLoRA Fine-Tuning of LLM Models} | Python, PyTorch
% \end{twocolentry}
% Concordia University – COMP433 Introduction to Deep Learning
% \vspace{0.10 cm}
% \begin{onecolentry}
%     \begin{highlights}
%     \item Fine-tuned GPT-Neo with QLoRA (4-bit NF4) for single-GPU training
%     \item Implemented LoRA adapters by modifying Q/K/V projections in PyTorch
%     \item Benchmarked accuracy, F1, and VRAM usage on SST-2, SQuAD v1.1, and AlpacaEval
%     \end{highlights}
% \end{onecolentry}


\section{Certifications}
\begin{twocolentry}{
    January 2026
}
Deep Learning Specialization – \textit{DeepLearning.AI / Coursera}
\end{twocolentry}
\begin{twocolentry}{
    April 2026
}
Anthropic Developer Certifications – \textit{Claude Code, Subagents, Claude Skills, MCP}
\end{twocolentry}


\vspace{-0.2cm}

  \section{Technical Skills}
  \begin{onecolentry}
      \textbf{Languages:} Python, Java, Bash, PowerShell, SQL \\
      \textbf{AI/ML:} PyTorch, Hugging Face, LoRA/QLoRA, LLM fine-tuning \\
      \textbf{Testing \& Observability:} JMeter, Selenium, Grafana, JUnit \\
      \textbf{Cloud \& DevOps:} AWS, Docker, Kubernetes, Helm, Jenkins, JFrog, GitLab CI \\
      \textbf{Databases:} MySQL, MongoDB, Cassandra, Neo4j, Redis
  \end{onecolentry}

\end{document}